\begin{textsample}{POS Dim 4 – 2015\_c5 – Score 1.00 – 2015\_c5/enade\_2015\_des\_1b821ff87a.txt}  \label{ex:f4_pos_163}
CAPÍTULO 5 DISTRIBUIÇÃO DOS CONCEITOS 5.1 PANORAMA NACIONAL DA DISTRIBUIÇÃO DOS CONCEITOS O Conceito Enade é um indicador de qualidade que avalia o desempenho dos estudantes a partir dos resultados do Enade.
Os valores possíveis vão de 1 ( pior situação ) a 5 ( melhor situação ).
As Unidades de \textbf{Observação} com apenas um ou sem nenhum concluinte participante não obtêm o Conceito Enade, ficando Sem Conceito ( SC ).
A Tabela 5.1 apresenta a quantidade e distribuição de cursos de Design participantes do Enade/2015, por faixa de conceito e Grande Região.
Mais uma vez, ressaltando que a diferença entre os cursos tabulados neste capítulo e no capítulo 2 corresponde aos cursos sem conceito, em princípio, aqueles sem alunos concluintes que participassem da prova ou que tivessem tão somente um aluno concluinte18.
Observando -se os dados da Tabela 5.1, nota -se que, dos 178 cursos participantes, 75 ( 42,1 \% ) classificaram -se com conceito 3, o valor modal.
Este foi também o conceito modal em mais três regiões : Nordeste ( 40,0 \% ), Sudeste ( 38,6 \% ) e Sul ( 50,7 \% ).
Na região Centro-Oeste, a moda foi o conceito 1 ( 37,5 \% ).
O conceito 2 foi o segundo mais frequente em nível nacional ( 27,0 \%, correspondendo a 48 cursos ), e o conceito 4, o terceiro ( 15,2 \%, correspondendo a 27 cursos ).
Houve, ainda, oito cursos ( 4,5 \% ) que receberam conceito 5 e nove ( 5,1 \% ) cursos que receberam conceito 1.
Dos 178 cursos de Design, 11 cursos ( 6,2 \% ) ficaram sem conceito ( SC ). 18 Estes cursos com somente um concluinte não foram considerados no capítulo 2.
A região Norte participou com cinco cursos ou 2,8 \% do total nacional.
Nesta região, o conceito 2 foi atribuído a três cursos dos cinco participantes, o que equivale a 60,0 \% do total regional, além de ser o conceito modal na região.
Nenhum curso ficou sem conceito ( SC ) ou com conceitos 1 ou 5.
Um curso foi avaliado com conceito 3 ( 20,0 \% ) e um curso obteve conceito 4 ( 20,0 \% ).
A região Nordeste participou com 20 cursos ou 11,2 \% do total nacional.
Nesta região, oito cursos, 40,0 \% em termos regionais, obtiveram conceito 3, o conceito modal para a região.
O conceito 2 foi obtido por cinco cursos ( 25 \% ), o conceito 4 por quatro cursos ( 20,0 \% ) e o 5 por apenas um curso ( 5,0 \% ).
Dois cursos obtiveram o conceito 1 ( 10,0 \% ) e nenhum curso ficou sem conceito ( SC ) nessa região.
Dos 70 cursos participantes da região Sudeste ( 39,3 \% do total nacional ), o conceito 3 também foi o conceito modal, e foi obtido por 27 cursos ( 38,6 \% ).
O conceito 1 foi atribuído a três cursos ( 4,3 \% ), e os conceitos 2 e 4, respectivamente, a 23 e dez cursos, correspondendo a 32,9 \% e 14,3 \%.
O conceito 5 foi atribuído a três cursos ( 4,3 \% ).
Apenas quatro cursos ( 5,7 \% ) ficaram sem conceito ( SC ).
Os 75 cursos da região Sul ficaram distribuídos em todas as faixas de conceitos.
Há predominância do conceito 3 com 50,7 \% ( conceito modal ), ou seja, 38 dos 75 cursos participantes na região Sul.
O conceito 1 foi atribuído a apenas um curso ( 1,3 \% ), o conceito 2, a 16 cursos ( 21,3 \% ) e o conceito 4 a 12 cursos ( 16,0 \% ).
O conceito 5 foi recebido por dois cursos ( 2,7 \% ).
Nessa região, seis cursos ( 8,0 \% ) ficaram sem conceito ( SC ).
Dos oito cursos participantes na região Centro-Oeste, três ( 37,5 \% ) receberam conceito 1, o conceito modal.
Os conceitos 2 e 3 foram atribuídos a um curso cada ( 12,5 \% cada ), e dois cursos ( 25,0 \% ) foram avaliados com conceito 5.
Nesta região, apenas um curso ficou sem conceito ( SC ), enquanto nenhum ficou com conceito 4.
O Gráfico 5.1 apresenta a distribuição do Conceito Enade segundo Grande Região : poligonais mais à esquerda representam Grandes Regiões com pior distribuição, e poligonais mais à direita, Grandes Regiões com melhores conceitos.
As regiões Norte ( linha azul ) e Sul ( linha roxa ) apresentam os cursos em IES com conceitos mais baixos.
A região Centro-Oeste ( linha preta ) apresenta um comportamento ambíguo, bimodal, com concentração nos melhores e piores conceitos e cruza as demais poligonais.
Cursos em IES das regiões Nordeste ( linha verde ) e Sudeste ( linha vermelha ) apresentam uma situação intermediária. 5.2 CONCEITOS POR CATEGORIA ADMINISTRATIVA E POR GRANDE REGIÃO A Tabela 5.2 apresenta a distribuição dos cursos participantes do Enade/2015 de Design, por Categoria Administrativa, de acordo com os conceitos por eles alcançados, segundo as Grandes Regiões brasileiras.
Dos 178 cursos participantes, 54 ( 30,3 \% ) eram ministrados em Instituições Públicas e 124 ( 69,7 \% ), em Instituições Privadas.
De acordo com as informações da Tabela 5.2, em termos nacionais, entre Instituições Públicas e Privadas nota -se que das oito IES que receberam o conceito 5, sete eram Públicas e uma Privada.
Dos 54 cursos participantes de IES Públicas, o conceito 4 foi o valor modal, atribuído a 25 cursos ( 46,3 \% ).
Entre os demais cursos de Instituições Públicas participantes, nenhum curso foi avaliado com conceito 1, seis cursos ( 11,1 \% da categoria ) obtiveram conceito 2, 14 cursos ( 25,9 \% ) receberam conceito 3 e, como já comentado, sete ( 13,0 \% ) foram alocados ao conceito 5.
Nesta categoria, dois cursos ( 3,7 \% ) ficaram sem conceito ( SC ).
Na rede privada, o conceito modal foi 3, com 61 cursos ( 49,2 \% ) dos 124 da categoria.
Entre os demais cursos participantes, nove ( 7,3 \% ) receberam o conceito 1 e 42 ( 33,9 \% ) o conceito 2.
O conceito 4 foi atribuído a dois cursos ( 1,6 \% ) e o conceito 5, a um curso ( 0,8 \% ).
Nesta Categoria Administrativa, nove cursos ( 7,3 \% ) ficaram sem conceito ( SC ).
Na análise por região, observa -se que, na região Norte, as Instituições Privadas participaram com quatro cursos ( 80,0 \% do total regional ), dos quais nenhum ficou sem conceito ( SC ) ou com conceitos 1, 4 ou 5.
O conceito modal para as Instituições Privadas na região foi 2, com três cursos, correspondendo a 75,0 \%, e apenas um curso foi alocado no conceito 3 ( 25,0 \% ).
As Instituições Públicas participaram com um curso na região Norte ( 20,0 \% do total regional ), o qual obteve conceito 4.
Na região Nordeste, a rede privada concentrou nove dos 20 cursos participantes, o equivalente a 45,0 \% do total da região.
Dentre os cursos oferecidos por IES Privadas no Nordeste, quatro cursos ficaram alocados no conceito 2, o conceito modal, três cursos obtiveram conceito 3 e dois cursos ficaram com conceito 1.
Nenhum curso recebeu conceito 4 ou 5 ou ficou sem conceito ( SC ).
As Instituições Públicas dessa região participaram com 11 cursos ( 55,0 \% do total da região ), dos quais cinco ( 45,5 \% do total dessa categoria ) obtiveram conceito 3, o valor modal, seguido por quatro cursos que receberam conceito 4 ( 36,4 \% ) e um curso ( 9,1 \% ) que foi avaliado com conceito 5.
Apenas um curso ( 9,1 \% ) recebeu conceito 2 e nenhum curso teve conceito 1 ou ficou sem conceito ( SC ) nesta combinação de categoria e região.
Na região Sudeste, a proporção de cursos da rede privada foi de 77,1 \%, correspondendo a 54 dos 70 cursos participantes.
Nesta categoria, na região Sudeste, o conceito modal foi 3 ( 24 cursos ).
Nesta combinação de Categoria Administrativa e Grande Região, quatro cursos ficaram sem conceito ( SC ).
Os demais foram avaliados com conceito 1 ( três cursos ), conceito 2 ( 21 cursos ), conceitos 4 e 5 ( um curso cada ).
Entre os 16 cursos oferecidos por Instituições Públicas na região Sudeste, o conceito modal foi 4 ( nove cursos ), e os demais cursos receberam os conceitos 2 ( dois cursos ), 3 ( três cursos ), 5 ( dois cursos ).
Na região Sudeste, nesta categoria, nenhum curso ficou sem conceito ( SC ).
As Instituições Privadas concentraram 51 dos 75 cursos participantes da região Sul, 68,0 \% do total regional.
Desses, 32 ficaram com conceito 3, o conceito modal, e 13 obtiveram conceito 2.
Nesta combinação de Categoria Administrativa e Grande Região, apenas um curso ( 2,0 \% ) ficou com conceito 1 e outro com conceito 4.
Nenhum curso ficou com o conceito 5 e quatro cursos ficaram sem conceito ( SC ).
As Instituições Públicas na região Sul participaram com 24 cursos ( 32,0 \% ) e o conceito modal foi 4, atribuído a 11 cursos ( 45,8 \% dos cursos dessa categoria, na região ).
Três cursos ( 12,5 \% ) foram avaliados com conceito 2, seis cursos ( 25,0 \% ) com conceito 3 e dois cursos ( 8,3 \% ) com conceito 5.
Nenhum curso recebeu o conceito 1 e dois cursos ficaram sem conceito ( SC ).
Na região Centro-Oeste, seis dos oito cursos participantes eram de Instituições Privadas ( 75,0 \% em termos regionais ).
Destes, três ( 50,0 \% ) concentraram -se no conceito 1, conceito modal.
Dentre os demais, um curso ( 16,7 \% ) recebeu conceito 2, e um ( 16,7 \% ) recebeu conceito 3.
Apenas um curso ficou sem conceito ( SC ) e nenhum

% matched lemmas: observação
\end{textsample}
